% Part: many-valued-logics
% Chapter: sequent-calculus
% Section: proof-theoretic-notions

\documentclass[../../../include/open-logic-section]{subfiles}

\begin{document}

\olfileid{mvl}{seq}{prf}

\olsection{Gagasan Téori Bukti}

\begin{explain}
Kaya déné kita wis netepaké sawetara gagasan semantik kang wigati
(validitas, akibat semantis, lan bisa dicukupi), saiki kita
netepaké \emph{gagasan téori bukti} kang cocog. Gagasan iki
ora ditetepaké miturut dicukupi utawa orané !!{sentence}s
ing !!{structure}s, nanging miturut !!{derivability} utawa
!!{nonderivability} saka sekèn tartamtu. Kasunyatan yèn
gagasan-gagasan iki cocog iku panemuan wigati; iki isi
\emph{téoréma kasahihan} lan \emph{téoréma kelengkapan}.
\end{explain}

\begin{defn}[Téoréma]
!!^a{sentence}~$!A$ diarani \emph{téoréma} yèn ana
!!a{derivation} ing~$\Log{LK}$ kanggo sekèn
$\quad \Sequent !A$. Kita nulis $\Proves !A$ yèn $!A$
iku téoréma, lan $\Proves/ !A$ yèn ora.
\end{defn}

\begin{defn}[!!^{derivability}]
!!^a{sentence}~$!A$ \emph{!!{derivable} saka} himpunan
!!{sentence}s~$\Gamma$, yaiku $\Gamma \Proves !A$,
yèn lan mung yèn ana himpunan pérangan winates
$\Gamma_0 \subseteq \Gamma$ lan runtutan $\Gamma_0'$
saka !!{sentence}s ing~$\Gamma_0$ saéngga $\Log{LK}$
!!{derive}s $\Gamma_0' \Sequent !A$. Yèn $!A$ ora
!!{derivable} saka $\Gamma$, kita nulis
$\Gamma \Proves/ !A$.
\end{defn}

Amarga ana aturan kontraksi, pelemahan, lan pertukaran,
urutan lan cacah !!{sentence}s ing~$\Gamma_0'$ ora wigati:
yèn sekèn $\Gamma_0' \Sequent !A$ !!{derivable},
mula $\Gamma_0'' \Sequent !A$ uga mangkono kanggo
saben $\Gamma_0''$ kang ngemot !!{sentence}s kang padha
karo~$\Gamma_0'$. Umpamané, yèn $\Gamma_0 = \{!B, !C\}$,
$\Gamma_0' = \tuple{!B, !B, !C}$ lan $\Gamma_0'' =
\tuple{!C, !C, !B}$ padha-padha runtutan kang mung
ngemot !!{sentence}s ing~$\Gamma_0$.
Yèn sekèn kanthi salah sijiné !!{derivable},
mula sekèn nganggo sijiné uga, umpamané:
\begin{prooftree}
  \AxiomC{}
  \Deduce$!B, !B, !C \fCenter !A$
  \RightLabel{\LeftR{\Contraction}}
  \UnaryInf$!B, !C \fCenter !A$
  \RightLabel{\LeftR{\Exchange}}
  \UnaryInf$!C, !B \fCenter !A$
  \RightLabel{\LeftR{\Weakening}}
  \UnaryInf$!C, !C, !B \fCenter !A$
\end{prooftree}
Wiwit saiki, yèn $\Gamma_0$ iku himpunan winates saka
!!{sentence}s, notasi $\Gamma_0 \Sequent !A$ tegesé
sekèn sembarang kang antesèdené runtutan !!{sentence}s
ing~$\Gamma_0$. Yèn perlu, kita anggep aturan kontraksi,
pertukaran, lan pelemahan wis ditrapaké kanthi meneng-meneng.

\begin{defn}[Konsistènsi]
Himpunan ukara~$\Gamma$ diarani \emph{ora konsisten}
yèn lan mung yèn ana himpunan pérangan winates
$\Gamma_0 \subseteq \Gamma$ saéngga $\Log{LK}$
!!{derive}s $\Gamma_0 \Sequent \quad$.
Yèn $\Gamma$ ora mangkono, yaiku yèn kanggo saben
$\Gamma_0 \subseteq \Gamma$ kang winates, $\Log{LK}$
ora !!{derive} $\Gamma_0 \Sequent \quad$,
mula diarani \emph{konsisten}.
\end{defn}

\begin{prop}[Refleksivitas]
\ollabel{prop:reflexivity}
Yèn $!A \in \Gamma$, mula $\Gamma \Proves !A$.
\end{prop}

\begin{proof}
Sekèn wiwitan $!A \Sequent !A$ iku !!{derivable},
lan $\{!A\} \subseteq \Gamma$.
\end{proof}
  
\begin{prop}[Monotonisitas]
\ollabel{prop:monotonicity}
Yèn $\Gamma \subseteq \Delta$ lan $\Gamma \Proves !A$,
mula $\Delta \Proves !A$.
\end{prop}

\begin{proof}
Anggepen $\Gamma \Proves !A$, yaiku ana $\Gamma_0
\subseteq \Gamma$ kang winates saéngga
$\Gamma_0 \Sequent !A$ iku !!{derivable}.
Amarga $\Gamma \subseteq \Delta$, $\Gamma_0$ uga
himpunan pérangan winates saka~$\Delta$.
Mula !!{derivation} kanggo $\Gamma_0 \Sequent !A$
uga nuduhaké $\Delta \Proves !A$.
\end{proof}

\begin{prop}[Transitivitas]
\ollabel{prop:transitivity}
Yèn $\Gamma \Proves !A$ lan $\{!A\} \cup \Delta \Proves
!B$, mula $\Gamma \cup \Delta \Proves !B$.
\end{prop}

\begin{proof}
Yèn $\Gamma \Proves !A$, ana $\Gamma_0 \subseteq \Gamma$
kang winates lan !!a{derivation}~$\pi_0$ kanggo
$\Gamma_0 \Sequent !A$. Yèn $\{!A\} \cup \Delta \Proves !B$,
mula kanggo sawatara $\Delta_0 \subseteq \Delta$ kang winates,
ana !!a{derivation}~$\pi_1$ kanggo $!A, \Delta_0
\Sequent !B$. Delengen !!{derivation} iki:
\begin{prooftree}
  \AxiomC{}
  \RightLabel{$\pi_0$}
  \Deduce$\Gamma_0 \fCenter !A$
  \AxiomC{}
  \RightLabel{$\pi_1$}
  \Deduce$!A, \Delta_0 \fCenter !B$
  \RightLabel{\Cut}
  \BinaryInf$\Gamma_0, \Delta_0 \fCenter !B$
\end{prooftree}
Amarga $\Gamma_0 \cup \Delta_0 \subseteq \Gamma \cup \Delta$,
iki nuduhaké $\Gamma \cup \Delta \Proves !B$.
\end{proof}

Mligi, iki ateges yèn $\Gamma \Proves !A$ lan $!A
\Proves !B$, mula $\Gamma \Proves !B$.
Uga, yèn $!A_1, \dots, !A_n \Proves !B$
lan $\Gamma \Proves !A_i$ kanggo saben~$i$,
mula $\Gamma \Proves !B$.

\begin{prop}
\ollabel{prop:incons}
$\Gamma$ ora konsisten yèn lan mung yèn
$\Gamma \Proves {!A}$ kanggo saben ukara~$!A$.
\end{prop}

\begin{proof}
Gladhen.
\end{proof}

\begin{prob}
Buktèkna \olref[mvl][seq][prf]{prop:incons}
\end{prob}

\begin{prop}[Kompaknes]
  \ollabel{prop:proves-compact}
  \begin{enumerate}
  \item Yèn $\Gamma \Proves !A$, ana himpunan pérangan winates
    $\Gamma_0 \subseteq \Gamma$ saéngga
    $\Gamma_0 \Proves !A$.
  \item Yèn saben himpunan pérangan winates saka~$\Gamma$
    konsisten, mula $\Gamma$ konsisten.
  \end{enumerate}
\end{prop}

\begin{proof}
  \begin{enumerate}
    \item Yèn $\Gamma \Proves !A$, ana himpunan pérangan winates
      $\Gamma_0 \subseteq \Gamma$ saéngga sekèn
      $\Gamma_0 \Sequent !A$ nduwèni !!a{derivation}.
      Akibaté, $\Gamma_0 \Proves !A$.
    \item Yèn $\Gamma$ ora konsisten, ana himpunan pérangan
      winates~$\Gamma_0 \subseteq \Gamma$ saéngga $\Log{LK}$
      !!{derive}s $\Gamma_0 \Sequent \quad$.
      Nanging $\Gamma_0$ iku himpunan pérangan winates
      saka~$\Gamma$ kang ora konsisten.
  \end{enumerate}
\end{proof}

\end{document}
